\section{The reach of the obstructions}\label{sec:scope}

An obstruction theorem is useful only if its boundary is drawn exactly. This
section draws the boundary of each of the main results, then states the
statements through which the conjecture is reached and proves the
implications of the literature that join them.

\subsection{Limits of the divisor obstruction}

\begin{remark}[Markman's construction]\label{rem:markmanscope}
The objects through which Markman proves the Weil classes of the sixfolds of
trivial discriminant algebraic \cite{Mar25a,Mar25b} are not obtained by
restriction from an ambient projective space, and they are not
$K$-equivariant. \Cref{sec:escape} analyses exactly which hypotheses they
escape and why, and \Cref{prop:whichescaped} lists the three. No obstruction
of this paper bears on those objects, and none needs to: their conclusion is
a theorem. The candidates proposed for cases that remain open, in
\cite[Section~12]{Mar25b} and \cite[Example~11.2.7]{Mar25c}, are another
matter, and \Cref{thm:candidatefails} and \Cref{cor:markmanquarticfails}
decide them.
\end{remark}

\begin{remark}[Twisted objects and gerbes]\label{rem:twisted}
The hypothesis of \Cref{thm:divisor} is that $\cF$ is an object of
$D^{b}(\PP^{N-1})$, so a twisted sheaf on a gerbe over $\PP^{N-1}$, or an
object of a derived category of twisted sheaves, is not covered. What is
covered, and what the proof uses, is that the Chern character of the
restricted object lies in the subring of $H^{*}(\Av,\QQ)$ generated by
divisor classes. Any construction whose input has that property is subject to
the conclusion, whatever category it is phrased in.
\end{remark}

\begin{remark}[Other ambient varieties]\label{rem:otherambient}
Replacing $\PP^{N-1}$ by another ambient variety $P$ changes the conclusion
exactly when $H^{*}(P,\QQ)$ is not generated by divisor classes, since the
proof only uses $\ch(\cF)\in D^{*}(P)$ and the fact that $i^{*}$ carries
$D^{*}(P)$ into $D^{*}(\Av)$. For $P$ a Grassmannian, a flag variety, or any
variety whose cohomology is generated in degree two, nothing changes. For $P$
with exceptional Hodge classes of its own the argument does not apply, but
then the problem has been moved rather than solved.
\end{remark}

\subsection{Limits of the character obstruction}

\begin{remark}[Constructions that are not equivariant]\label{rem:nonequivariant}
\Cref{thm:character} constrains eigenclasses. A construction whose output is
not a $\Kx$-eigenclass is not constrained, and by \Cref{rem:rightcharacter}
such a construction must break the naturality of the $K$-action in an
essential way: it must depend on a choice that is not $K$-invariant, and its
output must have components in both extreme summands of the grading at once.
These obstructions therefore point to constructions of that kind.
\end{remark}

\begin{remark}[Abelian varieties that are not very general]\label{rem:special}
At special points of the moduli of Weil type the endomorphism algebra grows,
the N\'eron--Severi group grows with it, and the Weil line can meet the divisor
subring. This is the split case, where Weil classes are known to be algebraic
by elementary means, and it is consistent with everything here: our statements
about $\HW(A)\cap D^{2n}(A)$ assume $\NS(A)\otimes\QQ=\QQ\eta$, which fails
there. \Cref{thm:character} itself makes no such assumption and holds at every
point, but at special points a class can be a sum of eigenclasses for several
characters and the conclusion applies to each summand separately.
\end{remark}

\subsection{Limits of the semiregularity statement}

\begin{remark}[Objects of positive rank]\label{rem:positiverankscope}
\Cref{thm:rankzero} and \Cref{cor:kernelfamily} both require the Chern character to
lie on the Weil line, so that $z\mapsto z\cdot\ch(\cE)$ has no degree-two
term. An object of nonzero rank has $\rk(\cE)\,z$ in that degree, the map is
injective, and nothing here applies. Markman's sheaves are of this kind, and
\Cref{rem:familyuniform} records why the distinction is not incidental to his
argument.
\end{remark}

\begin{remark}[The weaker semiregularity hypothesis]\label{rem:weakerscope}
What \Cref{cor:kernelfamily} obstructs is the injectivity of $\sigma_{\cE}$ on the
whole of $\Ext^{2}(\cE,\cE)$. The variant proposed in
\cite[Question 11.4]{Mar25b}, injectivity on the image of an evaluation map,
asks less. What we prove about it is confined to three groups of results.
In \Cref{thm:weakfails}, \Cref{prop:evaluation}, \Cref{thm:linebundle} and
\Cref{cor:noline} it fails. In \Cref{thm:factor}, for a secant object of
positive rank, it is shown to be equivalent to two conditions on the factor.
In \Cref{thm:candidatefails} and \Cref{cor:localobstruction} it fails for
Markman's candidate in dimension eight and for every secant object whose
support has two transversal branches of codimension two through a point. We
do not decide the variant for any other object of positive rank, nor whether
an affirmative answer to the question would follow from it.
\end{remark}

\begin{remark}[Deformation of the class itself]\label{rem:notnondeform}
A kernel is not an obstruction class. \Cref{cor:kernelfamily} says that a sufficient
criterion for the Weil class to stay algebraic is unavailable at the objects
in question; it does not say the class fails to stay algebraic, and no
statement in this paper does. \Cref{q:residual} is asked in both
directions.
\end{remark}

\subsection{Limits of the hyperplane obstruction}

\begin{remark}[Families that are not linear systems on \texorpdfstring{$X$}{X}]%
\label{rem:notlinear}
\Cref{thm:hyperplane} is about smooth members of an ample linear system on
$X$, and about the maps $j^{*}$ and $j_{*}$ they carry. It does not obstruct
an argument that deforms $X$ itself, that passes to a different variety by a
correspondence, or that uses a family in which $X$ appears as a fibre and the
other fibres have the same dimension. In the last case the inductive parameter
does not decrease, which is a different difficulty and the one recorded in
\Cref{rem:dimensionbookkeeping}.
\end{remark}

\begin{remark}[Coniveau]\label{rem:coniveauscope}
A Hodge class of positive coniveau is supported on a proper closed
subvariety, and the reduction to that subvariety is legitimate and lowers the
dimension. Nothing here obstructs it. The obstruction concerns classes of
coniveau zero, which is where the primitive middle cohomology of a general
variety sits, and the generalised Hodge conjecture predicts that these are
precisely the ones a coniveau argument cannot reach.
\end{remark}

\begin{remark}[The Kuga--Satake route]\label{rem:ksscope}
For a weight-two Hodge structure of K3 type the Kuga--Satake construction
produces an abelian variety $\mathrm{KS}(T)$ and an injection of $T$, twisted
once, into the endomorphisms of its first cohomology, carrying the Hodge class
in question to a class of type $(0,0)$, that is to an endomorphism, whose graph
is a cycle. The conclusion is transported back by a correspondence, not by a
restriction, so \Cref{thm:hyperplane} does not apply; this is
\Cref{prop:ksnotrestriction}. The route is limited by the requirement
$h^{2,0}=1$, and \Cref{prop:normpart} locates that limit exactly for abelian
varieties of Weil type: no rational sub-Hodge structure of $H^{2}(A,\QQ)$ is of
K3 type unless $n=1$. \Cref{sec:ks} treats the construction in full, including
the algebraicity of the correspondence itself.
\end{remark}

\subsection{The Weil statements and the closure map}
\label{ssec:closuremap}

The results above close several routes and narrow one. The statements
through which the conjecture is reached form a list, and we record each
entry with its exact statement and with what this paper proves about it.
\Cref{prop:descent} records the one way in which the entries on Weil classes
are linked.

Fix an imaginary quadratic field $K=\QQ(\sqrt{-d})$, an integer $n\ge2$, and a
class $\delta\in\QQ^{\times}_{>0}/\Nm_{K/\QQ}(\Kx)$, the discriminant of the
hermitian form. Write
\begin{quote}
$\mathbf{W}(K,n,\delta)$: every abelian $2n$-fold of $(K,-1,n)$-Weil type with
discriminant $\delta$ has algebraic Weil classes.
\end{quote}

\begin{proposition}[The Weil statements, and how they are linked]\label{prop:separate}
For each triple $(K,n,\delta)$ the following hold.
\begin{enumerate}[label=\textup{(\roman*)},leftmargin=2.6em]
\item $\mathbf{W}(K,n,\delta)$ is equivalent to the algebraicity of the single
  class $\omega_{s}$ at every point $s$ of the period domain $\cD_{n,n}$, and
  that domain is connected; this is \Cref{cor:exactform}.
\item $\mathbf{W}(K,3,\delta_{0})$ holds for the trivial discriminant
  $\delta_{0}$ \cite{Mar25a,Mar25b}, and with it, by \Cref{prop:descent},
  $\mathbf{W}(K,2,\delta)$ for every $K$ and every $\delta$, which is how
  Markman obtains the fourfolds.
\item The period domains attached to distinct triples are disjoint, and
  \Cref{thm:residual} and \Cref{cor:exactform} are statements about one
  triple at a time. The triples are nevertheless not independent: by
  \Cref{prop:descent}, $\mathbf{W}(K,N,\delta'')$ implies
  $\mathbf{W}(K,n,\delta)$ for every $n<N$ and every $\delta$, so the families
  of trivial discriminant in infinitely many dimensions give every imaginary
  quadratic family, and the families of a CM field of degree at least four
  give those of the imaginary quadratic fields it contains.
\end{enumerate}
\end{proposition}

\begin{proof}
(i) is \Cref{cor:exactform} together with \Cref{prop:weildomain}. (ii) is the
cited work together with \Cref{prop:descent}(i); \Cref{rem:residualmap}
records that the construction of \Cref{ssec:split} reaches only $\delta_{0}$.
For (iii), the discriminant is a discrete invariant of the hermitian form and
is constant on $\cD_{n,n}$ by \Cref{prop:splitgeom}(ii), while $n$ and $K$ are
fixed by the endomorphism algebra, so the domains are disjoint; the links are
\Cref{prop:descent}, which passes between families through a product with a
Weil surface or through a scalar extension, and not through a deformation.
The product argument of Schoen that Markman uses to pass from sixfolds to
fourfolds is an instance of such a link.
\end{proof}

\begin{proposition}[Descent and scalar extension]\label{prop:descent}
Let $F$ be a CM field of degree $2m$ and write $\mathbf{W}(F,n,\delta)$ for
the statement that every abelian variety of $(F,n,\delta)$-Weil type has
algebraic Weil classes, so that $\mathbf{W}(K,n,\delta)$ is the statement
above when $m=1$.
\begin{enumerate}[label=\textup{(\roman*)},leftmargin=2.6em]
\item \emph{Descent.} Let $B$ be of $(F,n,\delta)$-Weil type and $Y$ of
  $(F,1,\delta')$-Weil type. Then $X=B\times Y$, with the diagonal action of
  $F$ and the product polarisation, is of $(F,n+1,\delta\delta')$-Weil type,
  and the correspondence
  \[
    P(x)=pr_{B*}\bigl(x\cdot pr_{Y}^{*}(\eta_{Y}^{2m-2}\,y')\bigr),
  \]
  with $\eta_{Y}$ the polarisation of $Y$ and $y'$ a nonzero Weil class of
  $Y$, maps the Weil classes of $X$ isomorphically onto those of $B$. Since
  $y'$ lies in $H^{2}(Y)$ it is algebraic, so $P$ is algebraic, and every
  discriminant, a totally positive class in
  $F_{0}^{\times}/\Nm_{F/F_{0}}(F^{\times})$, is that of some $Y$.
  Hence $\mathbf{W}(F,N,\delta'')$ for one $N$ and one $\delta''$ implies
  $\mathbf{W}(F,n,\delta)$ for every $n<N$ and every $\delta$.
\item \emph{Scalar extension.} Let $K\subset F$ be CM fields and $B$ of
  $(K,n,\delta)$-Weil type. Then $B_{F}=B\otimes_{\cO_{K}}\cO_{F}$ is of
  $(F,n,\iota(\delta))$-Weil type, and for $j\colon B\to B_{F}$,
  $b\mapsto b\otimes1$, one has
  $j^{*}\Tr_{F/\QQ}(a\det_{F})=\Tr_{K/\QQ}(\Tr_{F/K}(a)\det_{K})$, so $j^{*}$
  maps the Weil classes of $B_{F}$ onto those of $B$. Hence
  $\mathbf{W}(F,n,\iota(\delta))$ implies $\mathbf{W}(K,n,\delta)$; every
  imaginary quadratic $K$ lies in the quartic CM field $K(\sqrt2)$.
\end{enumerate}
\end{proposition}

\begin{proof}
(i) $H^{1}(X)=H^{1}(B)\oplus H^{1}(Y)$ as $F$-spaces and the hermitian form
of $X$ is the orthogonal sum, of signature $(n+1,n+1)$ at every real place;
its normalised discriminant is the product, and $H_{Y}=\mathrm{diag}(1,-t)$
with $t$ totally positive realises every class. The Weil classes of $X$ are
$\bigoplus_{\sigma}\bigwedge^{2n}V_{B,\sigma}\otimes\bigwedge^{2}V_{Y,\sigma}$,
because the top exterior power of a direct sum is the tensor product of the
top powers. On the $\sigma$-component, $P$ is multiplication by
$\int_{Y}w_{Y}^{\sigma}\,\eta_{Y}^{2m-2}\,y'^{\bar\sigma}$, which is nonzero
since $w_{Y}^{\sigma}\wedge y'^{\bar\sigma}$ fills the $(\sigma,\bar\sigma)$
block and $\eta_{Y}$ is nondegenerate on the others; the remaining pairings
vanish for weight reasons under the norm-one torus of $F$. So $P$ is
rational, diagonal over $\CC$ with nonzero entries, and algebraic, and
$\omega(B)$ is algebraic whenever $\omega(X)$ is; \Cref{lem:onecomponent}
passes from one Weil class to all of them.

(ii) $H^{1}(B_{F})=H^{1}(B)\otimes_{K}F$, each eigenspace of $F$ is an
eigenspace of $K$ of signature $(n,n)$, the hermitian form is $H_{B}\otimes F$ with determinant $\det H_{B}$, and a norm
from $K$ is a norm from $F$. The identity for $j^{*}$ is the multiplicativity
of the determinant under restriction of scalars, and $\Tr_{F/K}$ is onto.
\end{proof}

The case $K$ imaginary quadratic, $n=2$ and $N=3$ of (i) is Proposition~10
of Schoen's addendum \cite{Sch98}, the argument through which Markman obtains
every Weil fourfold, at every discriminant, from the sixfolds of trivial
discriminant \cite{Mar25a}. Schoen multiplies a Weil fourfold $A$ by a Weil
surface $A'$ of the complementary discriminant, identifies $A\times A'$ up to
isogeny with a member of a family of Weil sixfolds of fixed discriminant, and
pushes forward the cycles $z\cdot(A\times D)$, for $D$ a divisor on $A'$
whose class spans its Weil classes; this is the correspondence $P$ of (i)
with $m=1$, and it involves no degeneration. What (i) adds is that the same
argument works in every dimension and for every CM field. As a consequence
the triples $(K,n,\delta)$ are not separate problems. The trivial
discriminant families in infinitely many dimensions give every imaginary
quadratic family, and by (ii) the split families of the CM fields of degree
at least four give the imaginary quadratic fields as well. In particular the
base points of nontrivial discriminant supplied by \Cref{thm:everyfamily} are
not needed on any route to the conjecture, since the split families already
carry base points \eqref{eq:splitclosed}; what \Cref{thm:everyfamily} shows
is that propagation can be started in every family, not that it has to be.
\Cref{fig:descent} draws the web of families that \Cref{prop:descent} links.

\begin{figure}[tbp]
\centering
  \includegraphics[width=0.94\textwidth]{fig_descent}
\caption{The Weil families of an imaginary quadratic field $K$ and of a
quartic CM field $F\supset K$, one cell for each $n$ and each kind of
discriminant, grass where the Weil classes are proved algebraic
(divisor classes at $n=1$; Markman \cite{Mar25a,Mar25b} at $n=2$ and at
$n=3$, $\delta=\delta_{0}$) and clay where they rest on the open
propagation statement. The indigo arrows
are the descent of \Cref{prop:descent}(i), from the trivial discriminant in
dimension $n+1$ to every discriminant in dimension $n$; the dashed teal
arrows are the scalar extension of \Cref{prop:descent}(ii). The dashed frame
is the propagation statement for the split families of the field of higher
degree, which by \Cref{thm:closure}(ii) is the only form of (P2) that any
minimal route to the conjecture uses: every clay cell of the left panel lies
below it.}
\label{fig:descent}
\end{figure}

\noindent
The list is as follows.

\begin{enumerate}[label=\textup{(C\arabic*)},leftmargin=3.0em,itemsep=4pt]
\item \emph{One secant object on one fourfold.} There should exist a
  principally polarised abelian fourfold $(X,\Theta)$ and a perfect complex
  $\cF$ on it with $\ch(\cF)=u_{\Theta}+3v_{\Theta}$ satisfying
  \eqref{eq:heisenberg}. By \Cref{thm:factor} and \Cref{cor:markmancondition}
  the transform $\Phi(\cF\boxtimes\cF)$ would then satisfy the weakened
  criterion, and by \Cref{prop:markman8} the three conditions of Markman's
  strategy that concern only the Chern character hold at the point $(1,3)$
  for any such $\cF$, since they depend on $\ch(\cF)$ alone. Turning this
  into algebraicity needs one further input. The weakened criterion enters
  in place of the hypothesis Markman verifies, and the survival of the
  semiregularity theorem under the substitution is exactly
  \cite[Question 11.4]{Mar25b}. So (C1) is a demand on a construction, and an affirmative answer to
  that question is a separate prerequisite standing beside it;
  \Cref{thm:closure} places both. What this paper contributes is the exact
  form of the demand. The local part of the criterion is a condition on depth
  and nothing else: if the support is Cohen--Macaulay there is no local
  obstruction at any point (\Cref{thm:cmvanishing}); if it is not, the local
  $\Ext^{2}$ is nonzero at the bad points, and where two smooth branches of
  codimension two meet transversally the criterion fails outright
  (\Cref{thm:candidatefails}, \Cref{cor:localobstruction},
  \Cref{cor:cmdichotomy}). The candidate of
  \cite[Section 12]{Mar25b} is of the second kind, and it cannot be moved to
  the first: separating the branches is what would restore the depth, and
  \Cref{rem:normalisationcm} shows that doing so moves the Chern character
  off the secant plane. What is needed is therefore a support without such
  points. For a Cohen--Macaulay support the identity \eqref{eq:heisenberg}
  is the single vanishing of the image of $A_{\theta}A_{\theta'}$ in
  $H^{1}(Z,\cE xt^{1}(\cI_{Z},\cI_{Z}))$, and the sheaf must in addition
  extend to first order along the polarised deformations
  (\Cref{thm:factor}(iii)). Such a support has one of two
  shapes (\Cref{cor:disjointroutes}): a smooth one, where the invariants are forced
  outright and $d\le b^{2}$ for the twist $b\Theta$, so $d\le9$ for the twist
  $3\Theta$ of \Cref{setup:markman} (\Cref{thm:smoothinvariants},
  \Cref{cor:fourdiscriminants}), and a singular one. In neither may the
  Hilbert--Burch resolution split into powers of $\Theta$ or into simple
  semi-homogeneous bundles (\Cref{thm:nosplit}, \Cref{thm:noblocks}), and
  no complete intersection occurs (\Cref{thm:noci}).
\item \emph{The same for every $n$ and every discriminant.} A positive answer
  to (C1) reaches one triple, and the discriminant of that triple is the
  trivial one, because the object lives on $X\times\widehat X$ and every split
  member has trivial discriminant by \Cref{prop:splitgeom}(ii). The other
  triples are reached through \Cref{prop:descent}: a positive answer in
  infinitely many dimensions gives $\mathbf{W}(K,n,\delta)$ for every $n$ and
  every $\delta$, so only uniformity in $n$ is asked for, and (C1) at $n=4$,
  with Question~11.4, already gives every sixfold. No answer of this kind
  reaches the CM fields of degree at least four. What this paper
  contributes: three of the four ingredients of \Cref{q:higher} are supplied
  for every $n$ and every $K$ (\Cref{thm:twoofthefour}), including the
  existence of the secant sheaves themselves (\Cref{thm:secantexists}); the
  fourth is isolated exactly (\Cref{rem:deformation}).
\item \emph{CM fields of degree larger than two.} For a CM field $F$ of
  degree $2m>2$ acting on an abelian variety, the Weil classes span a space of
  dimension $2m$, and the secant framework of \Cref{sec:construction}, which
  is built on a plane meeting the spinor variety in a conjugate pair, does not
  apply to them as it stands. Markman has extended it to these fields
  \cite{Mar25c}: for an abelian variety $X$ with real multiplication by the
  totally real subfield $F_{0}$ of $F$, a polarisation
  $\Theta\in\bigwedge^{2}_{F_{0}}H^{1}(X,\QQ)$ and $q\in F_{0}$ with
  $F=F_{0}(\sqrt{-q})$, he makes $F$ act on $X\times\widehat{X}$ with split
  Weil type, replaces the plane by a rational subspace $B$ of dimension
  $2^{[F_{0}:\QQ]}$ of the half-spin representation that is secant to the
  spinor variety, and shows that for sheaves $F_{1},F_{2}$ on $X$ with
  characters in $B$ the class $\kappa(E)=\ch(E)\exp(-c_{1}(E)/\rk E)$ of
  $E=\Phi(F_{1}\boxtimes F_{2}^{\vee})$ stays of Hodge type on every member
  of the family, and that where it stays algebraic all Weil classes are
  algebraic \cite[Theorem~1.1.2, Proposition~10.2.1]{Mar25c}. He exhibits
  such a pair meeting his non-degeneracy criterion when $F$ is biquadratic
  and $X$ is the Jacobian of a curve of genus four with real multiplication
  by $F_{0}$ and an automorphism acting through a unit of norm one
  \cite[Example~11.2.7, Lemma~11.2.8]{Mar25c}. That $\kappa(E)$ stays
  algebraic would follow from a semiregularity property of $E$, which is not
  addressed there \cite[Question~11.2.2]{Mar25c}; for the pair of that
  example the weaker property asked in the question fails at the points of
  the glued curves (\Cref{cor:markmanquarticfails}). No result of this paper
  depends on a theorem of \cite{Mar25c}; \Cref{ssec:quartic} and
  \Cref{ssec:sextic} use its construction as their setting. What this paper
  contributes: everything except the propagation. By \Cref{sec:cmfields}
  each such family has a connected period domain, a flat space of Weil
  classes, a dense set of rational orbits, a base point at a CM member where
  the classes are written as polynomials in divisor classes
  (\Cref{thm:cmbasepoint}), and the two propagation criteria
  (\Cref{thm:cmpropagation}); so the higher CM fields are the same
  propagation problem as the quadratic ones, and not a separate entry. For a
  quartic field at $n=2$, the smallest case of higher degree,
  \Cref{ssec:quartic}
  determines how far divisor classes, scalar extension and correspondences
  reach, and computes the number that the polarised criterion asks of one
  complex. For a sextic field, \Cref{ssec:sextic} shows that the dimension
  count, which excludes a minimal Orlov product over a quartic field, does not
  exclude one over a sextic field, and gives the profile such an object would
  have. \Cref{app:albert} places the imaginary quadratic case inside Albert's
  classification.
\item \emph{The non-split sixfolds.} $\mathbf{W}(K,3,\delta)$ for
  $\delta\ne\delta_{0}$; the case $\delta=\delta_{0}$ is Markman's
  \cite{Mar25a}. By \Cref{prop:descent}(i) it would follow from
  $\mathbf{W}(K,4,\delta_{0})$, and by \Cref{prop:descent}(ii) from
  $\mathbf{W}(K(\sqrt2),3,\iota(\delta))$. What this paper contributes: a base
  point in every family, at a member with quaternionic multiplication
  (\Cref{thm:everyfamily}); the locus of such members is a Lagrangian
  Grassmannian of codimension $n(n-1)/2$ (\Cref{prop:lagrangian}), and
  the rest of the family can be reached from it only by propagation.
\item \emph{Exceptional classes that are not of Weil type.} On an abelian
  variety the Hodge ring can contain exceptional classes that the Weil
  construction does not produce. Through dimension five it does not
  \cite{MZ99}, and in dimension eight it does: the self-product of an abelian
  fourfold of Mumford type carries two Hodge classes that are neither
  products of divisor classes nor Weil classes of any quotient
  (\Cref{prop:mumfordproduct}). So the Hodge conjecture for abelian varieties
  does not reduce to Weil classes, and the statement needed is the
  algebraicity of these classes themselves. What this paper contributes: the
  example, made exact, and the forms of their algebraicity studied in
  \Cref{sec:beyond}.
\item \emph{Varieties that are not abelian.} A reduction of the Hodge
  conjecture for arbitrary smooth projective varieties to abelian varieties
  cannot be of Hodge-theoretic kind: the $H^{3}$ of a very
  general quintic threefold lies outside the tensor category generated by the
  cohomology of abelian varieties (\Cref{prop:noabeliantype}), so its Hodge
  structure cannot be realised inside that of an abelian variety by any
  correspondence. \Cref{thm:hyperplane} and \Cref{thm:returntrip} close the
  two routes through hyperplane sections that would be the natural candidates.
  What this paper contributes: those three obstructions and, on the positive
  side, special cases only (\Cref{prop:mumfordfivefold},
  \Cref{rem:f3primefrontier}). Nor is this entry the complement of the others. The varieties
  that are not abelian include $A\times\PP^{1}$ for every abelian variety
  $A$, and the
  conjecture for $A\times\PP^{1}$ gives it for $A$, so the conjecture for
  the varieties that are not abelian is the whole conjecture
  (\Cref{prop:f3ishc}).
\end{enumerate}

\subsection{Implications from the literature}\label{ssec:bullseye}

The theorems of this paper are not the only implications that end in the
conjecture. The literature supplies others, and a frontier drawn without
them overstates what is needed. This subsection proves those implications, or
quotes them where they are theorems of the literature, so that
\Cref{ssec:closure} runs over both. Nothing in it is new
except the assembly, but the assembly changes the answer.

\begin{definition}[Three statements about all varieties]\label{def:lmv}
Here a variety is smooth, projective and defined over $\CC$, except the base
$S$ in \textup{(V)}.
\begin{enumerate}[leftmargin=3.0em,itemsep=3pt]
\item[\textup{(L)}] $B(X)$ holds for every variety $X$ (\Cref{def:Bconj}).
\item[\textup{(M)}] Every Hodge class on every variety $X$ is
  \emph{motivated} in the sense of Andr\'e \cite[\S2.1]{And96b}: it equals
  $\mathrm{pr}_{*}(\alpha\cup{*}\beta)$ for some variety $Y$, algebraic
  classes $\alpha,\beta$ on $X\times Y$, the projection
  $\mathrm{pr}\colon X\times Y\to X$, and the Lefschetz involution ${*}$ of
  $H^{*}(X\times Y,\QQ)$ with respect to some ample class.
\item[\textup{(V)}] For every smooth connected complex variety $S$, every
  smooth projective morphism $f\colon\cX\to S$ and every global section $\xi$
  of $R^{2p}f_{*}\QQ$ over $S^{\mathrm{an}}$: if $\xi_{s}$ is algebraic at one
  point $s\in S$, then $\xi_{t}$ is algebraic at every point $t\in S$.
\end{enumerate}
The Lefschetz involution of a variety $Z$ of dimension $e$ with ample class
$L$ sends $L^{r}x$, for $x$ primitive of degree $j$ and $0\le r\le e-j$, to
$\pm L^{e-j-r}x$, with a sign that depends only on $j$ \cite{And96b}; the
sign plays no role below.
\end{definition}

The statement (V) is the variational Hodge conjecture for algebraic classes;
it asks nothing about Hodge classes, only that algebraicity, once present,
survives deformation.

\begin{lemma}[Kleiman]\label{lem:lefinvolution}
Let $Z$ be a variety of dimension $e$ with ample class $L$, and suppose that
$B(Z)$ holds. Then the K\"unneth projectors $\pi_{k}$ of $H^{*}(Z,\QQ)$, the
projectors onto the summands $L^{r}P^{j}$ of the Lefschetz decomposition, and
the Lefschetz involution ${*}$ are all induced by algebraic cycles on
$Z\times Z$.
\end{lemma}

\begin{proof}
Operators induced by algebraic cycles on $Z\times Z$ are closed under sums
and rational multiples, and under composition, because the composite of two
algebraic correspondences is again algebraic \cite{Kle68}. The operator $L$
is induced by the image of a divisor of class $L$ under the diagonal, and $\Lambda$ by $B(Z)$. With
$[L,\Lambda]=H$ and $H$ acting on $H^{k}$ by $k-e$ (the proof of
\Cref{lem:LambdaL}), $H=L\Lambda-\Lambda L$ is induced by an algebraic cycle,
and since the numbers $k-e$ for $0\le k\le 2e$ are distinct,
\[
  \pi_{k}=\prod_{i\ne k}\frac{H-(i-e)}{k-i}
\]
is as well. Put $\Omega=H^{2}+2H+4\Lambda L$. For $x$ primitive of degree
$j$ put $m=e-j\ge0$. Then $\Lambda x=0$ and $Hx=-mx$, and the relation
\begin{equation}\label{eq:lambdapower}
  \Lambda L^{s}x=s(m-s+1)\,L^{s-1}x\qquad(s\ge1)
\end{equation}
follows by induction from $\Lambda L=L\Lambda-H$: the case $s=1$ is
\Cref{lem:LambdaL}, and
$\Lambda L^{s}x=L\Lambda L^{s-1}x-HL^{s-1}x
=\bigl((s-1)(m-s+2)-(2s-2-m)\bigr)L^{s-1}x=s(m-s+1)L^{s-1}x$.
So $\Omega x=(m^{2}-2m+4m)x=m(m+2)x$, and $\Omega$ commutes with $L$, since
$[L,H]=-2L$ and $[L,\Lambda L]=HL$ give $[L,\Omega]=-2(LH+HL)-4L+4HL=0$;
hence $\Omega$ acts on $L^{r}P^{j}$ by $m(m+2)$. Inside $H^{k}$ the summands
are $L^{r}P^{k-2r}$, with $m=e-k+2r$, and $m\mapsto m(m+2)$ is injective on
$m\ge0$, so Lagrange interpolation in $\Omega$ composed with $\pi_{k}$ gives
the projector $p_{k,r}$ onto $L^{r}P^{k-2r}$ as an algebraic operator.
Finally let $x\in P^{j}$ and $0\le r\le m$. If $2r\le m$, then
$L^{m-2r}(L^{r}x)=L^{m-r}x$. If $2r>m$, put $t=2r-m$, so $1\le t\le r$;
by \eqref{eq:lambdapower},
\[
  \Lambda^{t}L^{r}x=\prod_{i=0}^{t-1}(r-i)(m-r+i+1)\;L^{m-r}x,
\]
and every factor is at least $1$, because $r-i\ge r-t+1\ge1$ and
$m-r+i+1\ge m-r+1\ge1$. So on $L^{r}P^{j}$ the involution is $\pm L^{m-2r}$
or a nonzero rational multiple of $\Lambda^{2r-m}$, and
${*}=\sum_{k,r}c_{k,r}T_{k,r}\,p_{k,r}$ with rational $c_{k,r}$ and
$T_{k,r}$ a power of $L$ or of $\Lambda$. Every term is algebraic.
\end{proof}

\begin{proposition}[The conjecture is (L) and (M)]\label{prop:lefmot}
The Hodge conjecture for every variety is equivalent to \textup{(L)} and
\textup{(M)} together.
\end{proposition}

\begin{proof}
Assume the conjecture. Then \textup{(L)} holds by \Cref{thm:equivalence}. An
algebraic class $\xi$ on $X$ of dimension $d$ is motivated: take $Y$ a point,
$\alpha=\xi$, and $\beta=\varepsilon h^{d}$ with $h$ the ample class and
$\varepsilon=\pm1$ the sign for which ${*}(h^{d})=\varepsilon$, which exists
because the unit $1\in H^{0}$ is primitive of degree $0$; then
$\alpha\cup{*}\beta=\xi$. Every Hodge class is algebraic, so \textup{(M)}
holds.

Conversely assume \textup{(L)} and \textup{(M)}, and let $\xi$ be a Hodge
class on $X$. By \textup{(M)}, $\xi=\mathrm{pr}_{*}(\alpha\cup{*}\beta)$ with
$\alpha,\beta$ algebraic on $X\times Y$. By \textup{(L)}, $B(X\times Y)$
holds, so ${*}$ is induced by an algebraic cycle
(\Cref{lem:lefinvolution}); hence ${*}\beta$, the product
$\alpha\cup{*}\beta$ and its push-forward $\xi$ are algebraic. This is
Andr\'e's observation that motivated classes are algebraic as soon as the
Lefschetz involutions are \cite[\S2.1]{And96b}.
\end{proof}

\begin{proposition}[Two sources of (V)]\label{prop:lefvhc}
The Hodge conjecture implies \textup{(V)}, and so does \textup{(L)}.
\end{proposition}

\begin{proof}
Let $f\colon\cX\to S$, $\xi$ and $s$ be as in \textup{(V)}, with $\xi_{s}$
algebraic. As $S$ is smooth and connected it is irreducible, so two affine
open subsets meet, and it suffices to treat $S$ affine, joining $s$ to an
arbitrary $t$ through a point of the intersection of an affine neighbourhood
of each. Then $\cX$ is quasi-projective and, by resolution of singularities,
has a smooth projective compactification $\overline{\cX}$.

Assume the conjecture. Write $n=2p$ and
$\rho\colon H^{n}(\overline{\cX},\QQ)\to H^{0}(S^{\mathrm{an}},R^{n}f_{*}\QQ)$
for restriction; it is surjective by the theorem of the fixed part
\cite[4.1.1(ii)]{Del71}. For $t\in S$ the restriction
$r_{t}\colon H^{n}(\overline{\cX},\QQ)\to H^{n}(\cX_{t},\QQ)$ is $\rho$
followed by evaluation at $t$, and evaluation is injective because $S$ is
connected; so $K=\ker r_{t}=\ker\rho$ does not depend on $t$. It is a
sub-Hodge structure, and the Hodge structure $H^{n}(\overline{\cX},\QQ)$ is
polarisable, hence semisimple, so $H^{n}(\overline{\cX},\QQ)=K\oplus K'$
with $K'$ a sub-Hodge structure. Choose $\eta\in K'$ with $\rho(\eta)=\xi$.
Then $r_{s}|_{K'}$ is an injective morphism of Hodge structures and
$r_{s}(\eta)=\xi_{s}$ is of type $(p,p)$, so every component of $\eta$ of
type other than $(p,p)$ maps to zero and is zero: $\eta$ is a Hodge class.
So $\xi_{t}=r_{t}(\eta)$ is a Hodge class on $\cX_{t}$ for every $t$, and
algebraic by the conjecture.

Assume \textup{(L)}. The algebraic class $\xi_{s}$ is motivated (proof of
\Cref{prop:lefmot}), so $\xi_{t}$ is motivated for every $t$ by Andr\'e's
deformation theorem \cite[Th\'eor\`eme~0.5]{And96b}, and a motivated class is
algebraic under \textup{(L)} (proof of \Cref{prop:lefmot}).
\end{proof}

\begin{theorem}[Deligne, Andr\'e; Abdulali, Andr\'e]\label{thm:vhcab}
\textup{(V)} implies the Hodge conjecture for every abelian variety, and so
does \textup{(L)}.
\end{theorem}

\begin{proof}
Milne calls a class \emph{accessible} if it lies in the least family of
rational classes that contains the classes of algebraic cycles, is stable
under pull-back by morphisms of varieties, and contains $t_{t}$ for every $t$
whenever $(t_{s})_{s\in S}$ is a global section of $R^{r}f_{*}\QQ$ for a
smooth projective family over a smooth connected base and contains $t_{s}$
for one $s$; and he proves, following Deligne and Andr\'e, that every Hodge
class on an abelian variety is accessible \cite[\S5.1, Theorem~3]{Mil20}.
The algebraic classes contain the classes of algebraic cycles and are stable
under pull-back, and under \textup{(V)} they have the third property. So they
contain every accessible class, hence every Hodge class on an abelian
variety; this is \cite[Remark~2]{Mil20}. The statement for \textup{(L)}
follows by \Cref{prop:lefvhc}, and is the theorem of Abdulali and Andr\'e in
the form of \cite[Theorem~4]{Mil20}.
\end{proof}

\begin{remark}[The three in the literature]\label{rem:lmvknown}
$B(X)$ is known for curves, where it is immediate,
and for abelian varieties \cite[\S2]{Kle68}, and \Cref{cor:mumfordlefschetz}
below adds the total space of every Mumford family. \textup{(M)} holds for the Hodge
classes on abelian varieties \cite[Th\'eor\`eme~0.6.2]{And96b}.
\textup{(V)} holds in degree $2$, where by the argument of
\Cref{prop:lefvhc} the class $\xi_{t}$ is a Hodge class and the Lefschetz
theorem on $(1,1)$ classes applies; the one general mechanism for it in
higher degree is the semiregularity theorem \cite{Blo72,BF03,Pri24}, which is the
mechanism of \Cref{thm:semiregclosure}. \Cref{thm:vhcab} says what is at
stake for this paper: the single statement \textup{(V)}, which the conjecture
itself implies, gives the Weil classes of every CM field and the classes
beyond the Weil lines at once, by a route that uses none of the theorems
proved here.
\end{remark}

\subsection{The Lefschetz standard conjecture for abelian schemes over curves}
\label{ssec:newlefschetz}

The statement $B(X)$ is known for abelian varieties $X$. This subsection
proves it for the total space
of an abelian scheme over a curve whose invariant cycles are algebraic, and
so for the total space of every Mumford family. The proof assembles the
operator $\Lambda$ from a relative part, which is a Pontryagin product and is
algebraic on every abelian scheme, and a part along the base, which is
algebraic exactly when the invariant cycles are.

\begin{proposition}[$\Lambda$ as a Pontryagin product]
\label{prop:pontryaginlambda}
Let $A$ be an abelian variety of dimension $g$, $\ell\in H^{2}(A,\QQ)$ the
class of an ample line bundle, $D=\int_{A}\ell^{g}/g!$ and
$\gamma=\ell^{g-1}/(g-1)!$. Let $\mu\colon A\times A\to A$ be the addition
and $x\star y=\mu_{*}(p_{1}^{*}x\cup p_{2}^{*}y)$ the Pontryagin product.
Then the operator $\Lambda$ of \Cref{def:Bconj} for the class $\ell$ is
\[
  \Lambda(x)=D^{-1}\,x\star\gamma .
\]
In particular $\Lambda$ is induced by the algebraic cycle
$D^{-1}(p_{1},\mu)_{*}(p_{2}^{*}Z)$ on $A\times A$, for any cycle $Z$ of
class $\gamma$.
\end{proposition}

\begin{proof}
Put $V=H_{1}(A,\QQ)$, so that $H^{*}(A,\QQ)=\bigwedge^{*}V^{\vee}$ under cup
product and $H_{*}(A,\QQ)=\bigwedge^{*}V$ under the Pontryagin product of
homology, and let $\mathrm{PD}(x)=x\cap[A]$. Cap product makes
$H_{*}$ a module over $H^{*}$, and for a torus it is contraction; so
$\mathrm{PD}(\ell\cup x)=\ell\cap\mathrm{PD}(x)$ is the contraction
$\iota_{\ell}$ applied to $\mathrm{PD}(x)$. Poincar\'e duality commutes
with $\mu_{*}$ and carries $p_{1}^{*}x\cup p_{2}^{*}y$ to the cross product
of $\mathrm{PD}(x)$ and $\mathrm{PD}(y)$, without sign when $y$ has even
degree, so $\mathrm{PD}(x\star\gamma)=\mathrm{PD}(x)\cdot\beta$ with
$\beta=\mathrm{PD}(\gamma)\in\bigwedge^{2}V$.

Choose a basis $x_{1},y_{1},\dots,x_{g},y_{g}$ of $V^{\vee}$ with
$\ell=\sum_{i}d_{i}\,x_{i}\wedge y_{i}$, $d_{i}>0$, oriented so that
$\int_{A}x_{1}\wedge y_{1}\wedge\dots\wedge x_{g}\wedge y_{g}=1$; then
$\int_{A}\ell^{g}=g!\,d_{1}\cdots d_{g}$, so $D=d_{1}\cdots d_{g}$. Let
$e_{i},f_{i}$ be the dual basis of $V$. Then
$\gamma=\sum_{j}\bigl(\prod_{i\ne j}d_{i}\bigr)\prod_{i\ne j}x_{i}\wedge y_{i}$
and $\beta=\varepsilon D\sum_{j}d_{j}^{-1}e_{j}\wedge f_{j}$ with
$\varepsilon=\pm1$ independent of $j$, since the permutations of the pairs
preserve the orientation and exchange the terms. The algebra $\bigwedge^{*}V$
is the graded tensor product of the algebras
$\bigwedge^{*}\langle e_{i},f_{i}\rangle$, and $\iota_{\ell}$ and
multiplication by $D^{-1}\beta$ are sums over $i$ of even operators on the
$i$-th factor, namely $d_{i}\,\iota_{x_{i}\wedge y_{i}}$ and
$\varepsilon d_{i}^{-1}e_{i}\wedge f_{i}$. Operators on different factors
commute, and on the $i$-th factor, with basis $1,e_{i},f_{i},e_{i}\wedge
f_{i}$ and $\iota_{x_{i}\wedge y_{i}}(e_{i}\wedge f_{i})=\sigma$ for a sign
$\sigma$ fixed by the convention for contraction, the commutator
$[\iota_{\ell},D^{-1}\beta]$ acts by $\varepsilon\sigma$ on $1$, by $0$ on
$e_{i}$ and $f_{i}$, and by $-\varepsilon\sigma$ on $e_{i}\wedge f_{i}$, that
is by $\varepsilon\sigma(1-k_{i})$ in degree $k_{i}$. Summing over $i$,
$[\iota_{\ell},D^{-1}\beta]=\varepsilon\sigma(g-k)$ on $\bigwedge^{k}V$. As
$\mathrm{PD}$ carries $H^{q}$ to $\bigwedge^{2g-q}V$, the operator
$\Lambda'(x)=D^{-1}x\star\gamma$ satisfies
$[L,\Lambda']=\varepsilon\sigma(q-g)=\varepsilon\sigma H$ on $H^{q}(A)$.

The sign is fixed by one value. For $x$ and $y$ of complementary degrees,
$x\star y\in H^{0}(A)$ is $\int_{A\times A}p_{1}^{*}x\cup p_{2}^{*}y\cup
\mu^{*}[\mathrm{pt}]$, and $\mu^{*}[\mathrm{pt}]$ is the class of the
antidiagonal $\{(a,-a)\}$, so $x\star y=\int_{A}x\cup[-1]^{*}y$; for
$y=\gamma$, of even degree, this is $\int_{A}x\cup\gamma$. Hence
$\Lambda'(\ell)=D^{-1}\int_{A}\ell\cup\gamma=g$, while
$\Lambda(\ell)=\Lambda(L\cdot1)=g$ by \Cref{lem:LambdaL}. If
$\varepsilon\sigma=-1$, then $-\Lambda'$ would be the operator of
\Cref{def:Bconj} and $\Lambda'(\ell)=-g$. So $\varepsilon\sigma=1$, and
$\Lambda'=\Lambda$ by the uniqueness in \Cref{def:Bconj}. Finally, the
correspondence $(p_{1},\mu)_{*}(p_{2}^{*}Z)$ acts by
$x\mapsto p_{2*}\bigl(p_{1}^{*}x\cup(p_{1},\mu)_{*}p_{2}^{*}[Z]\bigr)
=\mu_{*}(p_{1}^{*}x\cup p_{2}^{*}\gamma)$, by the projection formula.
\end{proof}

\begin{theorem}[The Lefschetz operator of an abelian scheme over a curve;
cf.\ Tankeev \cite{Tan03}]
\label{thm:lefschetzfamily}
Let $C$ be a smooth projective connected curve and $\pi\colon W\to C$ an
abelian scheme of relative dimension $g$ with $W$ projective. Suppose that
every class in $H^{*}(W_{c},\QQ)$ invariant under the monodromy of
$R^{*}\pi_{*}\QQ$ is the restriction of an algebraic class on $W$. Then there
is an ample class $L_{W}$ on $W$ whose operator $\Lambda$ is induced by an
algebraic cycle on $W\times W$: $B(W)$ holds.
\end{theorem}

\begin{proof}
\emph{Step 1: the decomposition.} Let $[n]\colon W\to W$ be multiplication
by $n$ over $C$. It acts on $R^{q}\pi_{*}\QQ=\bigwedge^{q}R^{1}\pi_{*}\QQ$ by
$n^{q}$. The Leray spectral sequence degenerates \cite[4.1.1(i)]{Del71}, and
$[2]^{*}$ preserves the Leray filtration and acts on its graded pieces
$H^{p}(C,R^{q}\pi_{*}\QQ)$ by $2^{q}$. So $H^{k}(W,\QQ)=\bigoplus_{q}
H^{k}_{(q)}$, where $H^{k}_{(q)}$ is the $2^{q}$-eigenspace of $[2]^{*}$ and
maps isomorphically onto $H^{k-q}(C,R^{q}\pi_{*}\QQ)$; only $k-q\in\{0,1,2\}$
occurs. The projectors
$P_{q}=\prod_{q'\ne q}([2]^{*}-2^{q'})/(2^{q}-2^{q'})$ are induced by
algebraic cycles, polynomials in the graph of $[2]$. As $[2]^{*}$ is a ring
homomorphism, $H_{(q)}\cup H_{(q')}\subset H_{(q+q')}$; as
$\int_{W}[2]^{*}z=2^{2g}\int_{W}z$, the integral vanishes on $H_{(q)}$ for
$q\ne2g$, so Poincar\'e duality pairs $H_{(q)}$ with $H_{(2g-q)}$ only; and
$[2]_{*}=2^{2g}([2]^{*})^{-1}$ acts on $H_{(q)}$ by $2^{2g-q}$.

\emph{Step 2: the ample class.} Let $\cL$ be a relatively ample line bundle,
replaced by $\cL\otimes[-1]^{*}\cL$ so that it is symmetric. Then $[-1]^{*}$
fixes $c_{1}(\cL)$ and acts on $H_{(q)}$ by $(-1)^{q}$, so $c_{1}(\cL)$ has no
component in $H^{2}_{(1)}$, while $H^{2}_{(0)}=H^{2}(C,\QQ)$ is spanned by the
class $f$ of a fibre. So $c_{1}(\cL)=\ell+af$ with $\ell\in H^{2}_{(2)}$
algebraic and $a\in\QQ$. Put $L_{W}=\ell+mf=c_{1}(\cL)+(m-a)f$ with
$m\in\QQ$ so large that $L_{W}$ is a positive rational multiple of the class
of an ample line bundle, which is possible because $\cL$ is relatively ample.
Write $L_{\mathrm{rel}}$ and $L_{C}$ for cup product with $\ell$ and with
$mf$, so that $L_{W}=L_{\mathrm{rel}}+L_{C}$; on
$H^{p}(C,R^{q}\pi_{*}\QQ)$ the first raises $q$ by two and the second raises
$p$ by two.

\emph{Step 3: the relative operator.} Put $\gamma=\ell^{g-1}/(g-1)!$, which is
algebraic and lies in $H_{(2g-2)}$, and $D=\int_{W_{c}}\ell^{g}/g!$. Let
$\mu,\mathrm{pr}_{1},\mathrm{pr}_{2}\colon W\times_{C}W\to W$ be the addition
and the projections, and $T(x)=\mu_{*}(\mathrm{pr}_{1}^{*}x\cup
\mathrm{pr}_{2}^{*}\gamma)$. As in \Cref{prop:pontryaginlambda}, $T$ is
induced by the algebraic cycle $(\mathrm{pr}_{1},\mu)_{*}(\mathrm{pr}_{2}^{*}Z)$
on $W\times W$. Since $\pi\circ\mathrm{pr}_{1}=\pi\circ\mu$, the projection
formula gives $T(x\cup\pi^{*}b)=T(x)\cup\pi^{*}b$, so $T$ commutes with
$L_{C}$. Since $[2]\circ\mu=\mu\circ([2]\times_{C}[2])$,
$[2]_{*}T(x)=\mu_{*}(\mathrm{pr}_{1}^{*}[2]_{*}x\cup\mathrm{pr}_{2}^{*}
[2]_{*}\gamma)$, and for $x\in H_{(q)}$ this is
$2^{2g-q}\,2^{2}\,T(x)$; so $T$ maps $H_{(q)}$ into $H_{(q-2)}$. A relative
correspondence acts on the Leray spectral sequence through the maps it
induces on the local systems, so on $H^{p}(C,R^{q}\pi_{*}\QQ)$ the operator
$T$ is $H^{p}(C,T_{\bullet})$, with $T_{c}$ the Pontryagin product with
$\gamma|_{W_{c}}$ on the fibre. By \Cref{prop:pontryaginlambda}, $D^{-1}T_{c}$
is the operator $\Lambda$ of the fibre, so $\Lambda_{\mathrm{rel}}=D^{-1}T$
satisfies $[L_{\mathrm{rel}},\Lambda_{\mathrm{rel}}]=H_{\mathrm{rel}}$, where
$H_{\mathrm{rel}}$ acts on $H_{(q)}$ by $q-g$.

\emph{Step 4: the invariant classes.} Restriction maps $H^{q}(W,\QQ)$ onto the
monodromy invariants $I_{q}\subset H^{q}(W_{c},\QQ)$ \cite[4.1.1(ii)]{Del71},
and it kills $H^{q}_{(q')}$ for $q'<q$, which lies in positive Leray degree;
so it maps $H^{q}_{(q)}$ isomorphically onto $I_{q}$. Every element of
$H^{q}_{(q)}$ is algebraic: if $v\in I_{q}$ is the restriction of an algebraic
class $z$, then $P_{q}z$ is algebraic, lies in $H^{q}_{(q)}$ and restricts to
$v$. The subspace $I_{q}$ is a sub-Hodge structure, as the image of a
morphism of Hodge structures, and it is stable under cup product with
$\ell_{c}$ and under the operator $\Lambda$ of the fibre, which are the
restrictions of $L_{\mathrm{rel}}$ and $\Lambda_{\mathrm{rel}}$. The
polarisation of $H^{q}(W_{c},\QQ)$ attached to $\ell_{c}$ has the form
$(v,w)\mapsto\int_{W_{c}}v\cup S(w)$, with $S$ a polynomial in cup product
with $\ell_{c}$, the operator $\Lambda$ and the projectors of the fibre, and
its restriction to the sub-Hodge structure $I_{q}$ is again a polarisation
\cite{Voi02}. Since $S(I_{q})\subset I_{2g-q}$, the pairing
$I_{q}\times I_{2g-q}\to\QQ$ is nondegenerate.

\emph{Step 5: the operator along the base.} For $y\in H^{q}_{(q)}$ and
$w\in H^{2g-q}_{(2g-q)}$ one has $\int_{W}(y\cup f)\cup w=\int_{W_{c}}
y|_{W_{c}}\cup w|_{W_{c}}$. By Step 4 and Step 1 this shows that
$L_{C}\colon H^{q}_{(q)}\to H^{q+2}_{(q)}$ is injective, and it is onto
because $\dim H^{q+2}_{(q)}=\dim H^{2g-q}_{(2g-q)}=\dim I_{2g-q}=\dim I_{q}$,
by Poincar\'e duality on $W$ and on the invariants. Choose bases $(a_{j})$ of
$H^{q}_{(q)}$ and $(b_{k})$ of $H^{2g-q}_{(2g-q)}$, all algebraic by Step 4,
let $M_{jk}=\int_{W}a_{j}\cup mf\cup b_{k}$, which is invertible, and put
\[
  \Lambda_{C}(x)=\sum_{j,k}(M^{-1})_{kj}\Bigl(\int_{W}x\cup b_{k}\Bigr)a_{j}
  \qquad(x\in H^{q+2}(W,\QQ)),
\]
for every $q$. It is induced by the algebraic cycle
$\sum_{j,k}(M^{-1})_{kj}\,b_{k}\times a_{j}$. By Step 1, $\int_{W}x\cup b_{k}$
vanishes for $x\in H^{q+2}_{(q')}$ with $q'\ne q$, so $\Lambda_{C}$ is zero
on $H^{p}(C,R^{q}\pi_{*}\QQ)$ for $p\ne2$, and $\Lambda_{C}(L_{C}a_{i})=a_{i}$.
Hence $[L_{C},\Lambda_{C}]$ is $-1$ on $H_{(q)}^{q}$, $0$ on
$H_{(q)}^{q+1}$ and $+1$ on $H_{(q)}^{q+2}$: it is the operator $H_{C}$ that
acts on $H^{p}(C,R^{q}\pi_{*}\QQ)$ by $p-1$.

\emph{Step 6: the cross relations and the conclusion.} $[L_{C},
\Lambda_{\mathrm{rel}}]=0$ by Step 3. Both $L_{\mathrm{rel}}\Lambda_{C}$ and
$\Lambda_{C}L_{\mathrm{rel}}$ vanish off the summands with $p=2$, and on
$L_{C}y$ with $y\in H^{q}_{(q)}$ both give $\ell\cup y$, because
$\ell\cup y\in H^{q+2}_{(q+2)}$ and $L_{C}$ commutes with $L_{\mathrm{rel}}$;
so $[L_{\mathrm{rel}},\Lambda_{C}]=0$. Put
$\Lambda=\Lambda_{\mathrm{rel}}+\Lambda_{C}$, of degree $-2$. Then
\[
  [L_{W},\Lambda]=[L_{\mathrm{rel}},\Lambda_{\mathrm{rel}}]+[L_{C},\Lambda_{C}]
  =H_{\mathrm{rel}}+H_{C},
\]
which acts on $H^{p}(C,R^{q}\pi_{*}\QQ)\subset H^{p+q}(W)$ by
$(q-g)+(p-1)=(p+q)-\dim W$. So $\Lambda$ is the operator of \Cref{def:Bconj}
for $L_{W}$, by its uniqueness, and it is induced by an algebraic cycle by
Steps 3 and 5.
\end{proof}

\Cref{fig:lefschetzgrid} draws the summands of the proof in relative
dimension four.

\begin{corollary}[Mumford families]\label{cor:mumfordlefschetz}
Let $\pi\colon W\to C$ be a non-isotrivial abelian scheme of relative dimension
four over a smooth projective connected curve, $W$ projective, whose very
general fibre $X$ is a fourfold as in \Cref{prop:mumfordproduct}, with Hodge
group the group $G$ of Mumford's construction \cite{Mum69}, the restriction of
scalars from a totally real cubic field $F$ of the norm one group of a
quaternion algebra over $F$. Then $B(W)$ holds. Such families exist, over the
compact Shimura curves of \cite{Mum69}.
\end{corollary}

\begin{proof}
By Andr\'e's theorem \cite[Theorem~1]{And92} the connected monodromy group is
a normal subgroup of the derived Mumford--Tate group of $X$, which is $G$. It
is not trivial, because the family is not isotrivial, and $G$ is $\QQ$-simple,
since over $\overline{\QQ}$ its three factors are indexed by the embeddings of
$F$ and permuted transitively by the Galois group; so it is $G$. The
monodromy invariants in $H^{*}(X,\QQ)$ are therefore contained in the
invariants of $G$, which are the Hodge classes of $X$, and these are the
powers of the polarisation by \Cref{prop:mumfordproduct}(i). The powers
of $\ell$ restrict to the powers of the polarisation, so every invariant class
is the restriction of an algebraic class, and \Cref{thm:lefschetzfamily}
applies.
\end{proof}

\begin{figure}[tbp]
\centering
  \includegraphics[width=0.86\textwidth]{fig_lefschetzgrid}
\caption{The proof of \Cref{thm:lefschetzfamily} for relative dimension
four. The boxes are the summands $H^{p}(C,R^{q}\pi_{*}\QQ)$ of $H^{*}(W,\QQ)$,
split off by the eigenvalue $2^{q}$ of $[2]^{*}$ shown on the left; the dotted
diagonals join the summands of one $H^{k}(W,\QQ)$, $k=p+q$, and the dashed
magenta line is the axis $q=g$ of hard Lefschetz on the fibres. The relative
operators $L_{\mathrm{rel}}$ and $\Lambda_{\mathrm{rel}}$, in teal, move
within a column; $\Lambda_{\mathrm{rel}}$ is a Pontryagin product with an
algebraic class on every abelian scheme (\Cref{prop:pontryaginlambda}). The
base operators $L_{C}$ and $\Lambda_{C}$, in ochre, move along a row between
the outer columns and kill the middle one; $\Lambda_{C}$ is the displayed
cycle built from bases $(a_{j})$, $(b_{k})$ of the invariant summands, and is
algebraic exactly when the invariant cycles are. Below each column is the
weight of $H=[L_{W},\Lambda]$ on it, and the panel at the top is Step 6 of the
proof. The ochre boxes carry the invariant cycles and their images under
$L_{C}$; for a Mumford family they are the lines listed in magenta on the
right, spanned by $\ell^{i}f$ (and by $\ell^{i}$ in the left column), and the
other boxes of the outer columns are zero (\Cref{cor:mumfordlefschetz}).}
\label{fig:lefschetzgrid}
\end{figure}

\begin{remark}[Scope of the argument]\label{rem:newlefschetzscope}
The total space $W$ is a fivefold fibred in abelian fourfolds whose Hodge
groups are not defined by endomorphisms; it is not a product and not an
abelian variety. For abelian schemes over curves the mechanism of
\Cref{thm:lefschetzfamily} goes back to Tankeev \cite{Tan03}, who proves $B$
for the total space when certain canonically defined Hodge cycles are
algebraic; the argument here does not use his hypotheses, and we make no
claim of priority for \Cref{cor:mumfordlefschetz}. The difficulty lies in the hypothesis of
\Cref{thm:lefschetzfamily}.
Applied to the fibre square $W\times_{C}W\to C$ the theorem gives
$B(W\times_{C}W)$ as soon as every invariant class of $X\times X$ is the
restriction of an algebraic class on $W\times_{C}W$; the invariant classes
include the two exceptional classes of \Cref{prop:mumfordproduct}, so this
is the smallest case of (F2) in a global form. \Cref{thm:lefschetzfamily} does not bear on (P2), on the Kuga--Satake
class of \Cref{thm:mumfordks}, or on \textup{(L)} for varieties that are not
fibred in abelian varieties over a curve.
\end{remark}

\begin{proposition}[Powers of a Weil variety; cf.\ \cite{Mil22}]\label{prop:powers}
Let $X$ be of $(K,-1,n)$-Weil type with $n\ge2$ and $\Hg(X)=\SU(V,H)$. For
every $k\ge1$ the $\QQ$-algebra $\Hdg^{*}(X^{k})$ is generated by the divisor
classes of $X^{k}$ and the pull-backs $f_{a}^{*}\omega$ of the Weil classes
$\omega$ of $X$ along the homomorphisms $f_{a}=\sum_{j}a_{j}\,pr_{j}\colon
X^{k}\to X$, $a\in\ZZ^{k}$. Consequently, if the Weil classes of $X$ are
algebraic, the Hodge conjecture holds for every power of $X$. By
\cite{Mar25a} this is so for every such $X$ with $n=2$, and with $n=3$ and
trivial discriminant.
\end{proposition}

\begin{proof}
Write $H^{1}(X,\CC)=V_{+}\oplus V_{-}$ for the eigenspaces of $K$ and
$W=V_{+}$. The polarisation makes $V_{+}$ and $V_{-}$ isotropic and pairs them
perfectly, so $V_{-}=W^{\vee}$, and restriction to $V_{+}$ identifies
$\SU(V,H)_{\CC}$ with $\SL(W)$. As $\Hg(X^{k})$ is $\Hg(X)$ acting diagonally
on $H^{1}(X^{k},\CC)=(W\otimes U)\oplus(W^{\vee}\otimes U)$, $U=\CC^{k}$, the
Hodge classes of $X^{k}$, tensored with $\CC$, are the $\SL(W)$-invariants in
$\bigwedge^{*}\bigl((W\otimes U)\oplus(W^{\vee}\otimes U)\bigr)$. That algebra
is an $\SL(W)$-equivariant quotient of a tensor algebra, with an equivariant
section given by antisymmetrisation, so Weyl's first fundamental theorem for
$\SL(W)$, in its tensor form and applied before antisymmetrisation, shows
that its invariants are generated by the images of the contractions and the
determinant tensors: the classes
$D(u,u')=\sum_{i}(e_{i}\otimes u)\wedge(e_{i}^{\vee}\otimes u')$ and the
polarised determinants, which by polarisation are combinations of
$\det_{u}=\bigwedge_{i}(e_{i}\otimes u)$ and
$\det^{\vee}_{u}=\bigwedge_{i}(e_{i}^{\vee}\otimes u)$, $u,u'\in U$. For
$2n\ge3$ the $D(u,u')$ span the invariants of degree two, which are the
divisor classes tensored with $\CC$. With $\alpha_{\pm}$ spanning
$\bigwedge^{2n}V_{\pm}$ (\Cref{prop:weilline}) one has
$f_{a}^{*}\alpha_{+}=\det_{a}$ and $f_{a}^{*}\alpha_{-}=\det^{\vee}_{a}$, and
since $u\mapsto\det_{u}$ is polynomial and $\ZZ^{k}$ is Zariski dense in
$U$, these span the same space as all $\det_{u}$. So the $\QQ$-algebra
generated by $\NS(X^{k})_{\QQ}$ and the $f_{a}^{*}\omega$ is contained in
$\Hdg^{*}(X^{k})$ and has the same complexification. The generators are
algebraic when the Weil classes of $X$ are, and they are for $n=2$ and for
$n=3$ at trivial discriminant by \cite{Mar25a}.
\end{proof}

The statement for powers is known (compare \cite{Mil22}, and \cite{Li26}
for the powers of CM fourfolds); we include the short proof because the next
corollary uses it.

\begin{corollary}[Abelian schemes of small relative dimension]
\label{cor:lefschetzsmall}
Let $\pi\colon W\to C$ be an abelian scheme of relative dimension $g$ over a
smooth projective connected curve, with $W$ projective, and suppose that
every monodromy invariant class in $H^{*}(W_{c},\QQ)$ is a Hodge class, as
happens when the identity component of the Zariski closure of the monodromy
group is the Hodge group of a very general fibre.
\begin{enumerate}[label=\textup{(\roman*)},leftmargin=2.6em,itemsep=2pt]
\item If $g\le4$, then $B(W)$ holds.
\item If the fibres are of $(K,-1,n)$-Weil type with connected monodromy
  group $\SU(V,H)$, and $n=2$, or $n=3$ with trivial discriminant, then
  $B$ holds for every fibre power $W\times_{C}\dots\times_{C}W$. In
  particular $B(W_{C})$ of \Cref{cor:weilequiv} holds in these cases.
\end{enumerate}
\end{corollary}

\begin{proof}
By the theorem of the fixed part the invariant classes form a sub-Hodge
structure of every fibre, with a Hodge structure independent of the point, so
they are Hodge classes on every fibre.

(i) The Hodge conjecture holds for abelian varieties of dimension at most
four, by the Lefschetz theorem on $(1,1)$ classes and hard Lefschetz up to
dimension three and by \cite{Mar25a} in dimension four, so each invariant
class is algebraic on every fibre; \Cref{lem:spreading} makes it the
restriction of an algebraic class on $W$, and \Cref{thm:lefschetzfamily}
applies.

(ii) The fibre power is an abelian scheme with projective total space and
the monodromy acting diagonally. Its invariant classes are
$\SU(V,H)$-invariants, hence Hodge classes at every member with
$\Hg=\SU(V,H)$, an uncountable set of points, and there the Hodge conjecture
holds for the fibre by \Cref{prop:powers}; \Cref{lem:spreading} and
\Cref{thm:lefschetzfamily} conclude as before. For $W_{C}$ the connected
monodromy group is $\SU(V,H)$ when it is Zariski dense, which is the
hypothesis of \Cref{cor:weilequiv}(ii).
\end{proof}

The hypothesis on the monodromy in \Cref{cor:lefschetzsmall} can be checked
with Andr\'e's theorem \cite[Theorem~1]{And92}: the identity component of the
Zariski closure of the monodromy group is always a normal subgroup of the
derived Mumford--Tate group of a very general fibre, so the hypothesis holds
whenever the family is not isotrivial and the Hodge group of a very general
fibre is semisimple and $\QQ$-simple, as for the Mumford families of
\Cref{cor:mumfordlefschetz}.

\begin{proposition}[The total space of a Mumford family]
\label{prop:mumfordfivefold}
Let $W\to C$ be as in \Cref{cor:mumfordlefschetz}, with $C$ the compact
Shimura curve of \cite{Mum69} of torsion-free level, or a finite \'etale cover
of it. Then the Hodge conjecture holds for the fivefold $W$: its Hodge classes
of degree $2p$ are spanned by $\ell^{p}$ and $f\ell^{p-1}$, with $f$ the class
of a fibre.
\end{proposition}

\begin{proof}
By Step 1 of the proof of \Cref{thm:lefschetzfamily},
$H^{k}(W,\QQ)=\bigoplus_{q}H^{k}_{(q)}$ with $H^{k}_{(q)}\cong
H^{k-q}(C,R^{q}\pi_{*}\QQ)$, the summands being sub-Hodge structures cut out by
algebraic projectors. The summand with $k=q$ is the space of invariants, which
by the proof of \Cref{cor:mumfordlefschetz} is spanned by the powers of the
polarisation, restrictions of the powers of $\ell$. The summand with $k=q+2$
is $f\cup H^{q}_{(q)}$ by Step 5, spanned by the $f\ell^{q/2}$. It remains to
see that $H^{1}(C,R^{q}\pi_{*}\QQ)$ carries no Hodge class. For $q$ even its
weight $q+1$ is odd. For $q$ odd, over $\CC$ the local system splits as a sum
of $S^{k_{1}}V_{1}\otimes U_{\lambda}\otimes M_{\lambda}$, where $V_{1}$ is the
uniformising variation of the factor of $G$ split at the real place, $U_{\lambda}$ is
a unitary variation of type $(0,0)$, $M_{\lambda}$ is constant of type $(a,a)$
with $2a+k_{1}=q$, and $k_{1}\equiv q\pmod2$ because the central element
$(-1,1,1)$ acts on $V$ by $-1$. When $C$ is \'etale over the Shimura curve the
Kodaira--Spencer map of $V_{1}$ is an isomorphism, and then
$H^{1}(C,S^{k_{1}}V_{1}\otimes U_{\lambda})$ has only the types $(k_{1}+1,0)$
and $(0,k_{1}+1)$ \cite{Zuc79}; after the twist by $M_{\lambda}$ these are
never balanced, since $k_{1}\ge1$.
\end{proof}

Statements of the kind of \Cref{prop:mumfordfivefold} are studied for Kuga
fibre varieties, and we make no claim of priority.
The \'etale hypothesis is used only at the end of the proof, and it cannot be
dropped from the argument: over a ramified cover the Kodaira--Spencer map of
$V_{1}$ vanishes at the branch points, the Higgs field of $S^{k_{1}}V_{1}$ is
no longer an isomorphism, and the torsion of its cokernel contributes to
$H^{1}(C,R^{q}\pi_{*}\QQ)$ in balanced types.

\subsection{The two targets as cases of (L)}
\label{ssec:targets}

The argument of \Cref{ssec:newlefschetz} does more than add a case of
\textup{(L)}. For an abelian scheme over a curve it shows that \textup{(L)}
for the total space is exactly the statement that algebraicity propagates
along the curve. The base points of this paper then turn each of the two
targets of \Cref{rem:bullseye}, one Weil family and the first
classes beyond the Weil lines, into \textup{(L)} for one explicit variety. The
reductions are equivalences, so they locate the difficulty rather than remove
it. Along the way the Hodge conjecture is proved for the self-product of a
Mumford fourfold at every CM point (\Cref{thm:mumfordcm}), and the two
branches of \Cref{rem:twobranches} are shown not to be realisable by two
objects (\Cref{thm:nosum}), an implication whose hypothesis
\Cref{thm:p2false} shows is never met.

\begin{lemma}[Spreading along a curve]\label{lem:spreading}
Let $\pi\colon W\to C$ be a smooth projective morphism onto a smooth projective
connected curve, $W$ projective, and $\xi$ a global section of
$R^{2p}\pi_{*}\QQ$. If $\xi_{t}$ is algebraic for uncountably many $t\in C$,
there are an algebraic class $\tilde\xi$ on $W$ and an integer $d\ge1$ with
$\tilde\xi|_{W_{t}}=d\,\xi_{t}$ for every $t\in C$.
\end{lemma}

\begin{proof}
The proof of \Cref{thm:closedlocus} applies verbatim with $(\cA,\cS,\omega)$
replaced by $(W,C,\xi)$: the set of $t$ at which $\xi_{t}$ is algebraic is
the union, over countably many tuples of Hilbert polynomials and rational
coefficients, of the images of unions $T$ of connected components of fibre
products of relative Hilbert schemes, projective over $C$, on which
$\sum_{i}q_{i}[Z_{i}]=\xi$. Uncountably many $t$ force one image to be
infinite, hence equal to $C$. Choose an irreducible component of that $T$
surjecting onto $C$, two of its points over distinct points of $C$, and an
irreducible curve through them; its normalisation $C'$ maps finitely onto $C$,
of some degree $d$. Pulling the universal subschemes back to
$W'=W\times_{C}C'$ gives a cycle $z=\sum_{i}q_{i}[\cZ_{i}]$ whose restriction
to $W'_{c'}=W_{g(c')}$ is $\xi_{g(c')}$ for every $c'\in C'$. The projection
$h\colon W'\to W$ is finite and flat of degree $d$, and by base change
$h_{*}[z]$ restricts on $W_{t}$ to the sum of the restrictions over the
points of $C'$ above $t$, counted with multiplicity, which is $d\,\xi_{t}$.
\end{proof}

\begin{theorem}[\textup{(L)} for an abelian scheme over a curve is
propagation]\label{thm:lefschetzpropagation}
Let $\pi\colon W\to C$ be as in \Cref{thm:lefschetzfamily}, but with no
hypothesis on its invariant cycles, and keep the notation of its proof. Let
$\Lambda_{C}$ be the operator that is zero on $H^{p}(C,R^{q}\pi_{*}\QQ)$ for
$p\ne2$ and inverts $L_{C}\colon H^{q}_{(q)}\to H^{q+2}_{(q)}$.
\begin{enumerate}[label=\textup{(\roman*)},leftmargin=2.6em,itemsep=2pt]
\item The operator $\Lambda$ of \Cref{def:Bconj} for $L_{W}$ is
  $\Lambda_{\mathrm{rel}}+\Lambda_{C}$. So $B(W)$ holds for $L_{W}$ if and
  only if $\Lambda_{C}$ is induced by an algebraic cycle.
\item If $B(W)$ holds for $L_{W}$, then for every $c\in C$ and every algebraic
  class $y$ on $W_{c}$ the class $\Lambda_{C}(i_{c*}y)$ is algebraic; up to
  the factor $m$ it lies in $H^{0}(C,R^{q}\pi_{*}\QQ)$ and restricts on
  $W_{c}$ to the unique monodromy invariant class $v$ with
  $\int_{W_{c}}v\cup w=\int_{W_{c}}y\cup w$ for every invariant $w$ of
  complementary degree. In particular every global section of
  $R^{2p}\pi_{*}\QQ$ that is algebraic at one point is the restriction of an
  algebraic class on $W$, and is algebraic at every point.
\item If at one point $c_{0}$ every monodromy invariant class of
  $H^{*}(W_{c_{0}},\QQ)$ is algebraic, then $B(W)$ holds for $L_{W}$ if and
  only if every monodromy invariant class is the restriction of an algebraic
  class on $W$.
\end{enumerate}
\end{theorem}

\begin{proof}
(i) Steps 4 and 5 of the proof of \Cref{thm:lefschetzfamily} show that
$L_{C}\colon H^{q}_{(q)}\to H^{q+2}_{(q)}$ is an isomorphism without using
the hypothesis on invariant cycles: the nondegeneracy of
$I_{q}\times I_{2g-q}\to\QQ$ comes from the Hodge--Riemann relations alone.
So $\Lambda_{C}$ is defined, and Steps 5 and 6, which use only that
$\Lambda_{C}$ vanishes off $p=2$ and inverts $L_{C}$, give
$[L_{W},\Lambda_{\mathrm{rel}}+\Lambda_{C}]=H$. By the uniqueness in
\Cref{def:Bconj} this is $\Lambda$. As $\Lambda_{\mathrm{rel}}$ is induced by
an algebraic cycle (Step 3), $\Lambda$ is if and only if $\Lambda_{C}$ is.

(ii) The Gysin map $i_{c*}$ commutes with $[2]_{*}$, which acts on
$H^{q}(W_{c})$ by $2^{2g-q}$, so $i_{c*}y\in H^{q+2}_{(q)}$, the summand with
$p=2$. If $y$ is algebraic so is $i_{c*}y$, and by (i) so is
$\Lambda_{C}(i_{c*}y)$. Put $u=m\,\Lambda_{C}(i_{c*}y)\in H^{q}_{(q)}$, so
$f\cup u=i_{c*}y$. For $w\in H^{2g-q}_{(2g-q)}$ the projection formula gives
$\int_{W_{c}}u|_{W_{c}}\cup w|_{W_{c}}=\int_{W}(f\cup u)\cup w
=\int_{W}i_{c*}y\cup w=\int_{W_{c}}y\cup w|_{W_{c}}$, and $w|_{W_{c}}$ runs
over the invariants of complementary degree (Step 4); by the nondegeneracy of
Step 4 this determines $u|_{W_{c}}$. If $y=\xi_{c}$ for a global section
$\xi$, then $v=\xi_{c}$ has the property, so $u|_{W_{c}}=\xi_{c}$, and $u$,
being a global section of $R^{q}\pi_{*}\QQ$ equal to $\xi$ at $c$, restricts
to $\xi_{t}$ at every $t$.

(iii) If every invariant class is a restriction of an algebraic class, $B(W)$
holds by \Cref{thm:lefschetzfamily}. Conversely, under $B(W)$, (ii) applied at
$c_{0}$ to each invariant class makes it the restriction of an algebraic
class.
\end{proof}

\Cref{fig:propagation} draws part (ii) along a compact Shimura curve.

\begin{figure}[tbp]
\centering
  \includegraphics[width=0.84\textwidth]{fig_propagation}
\caption{\Cref{thm:lefschetzpropagation}(ii) along a compact Shimura curve
$C$. An algebraic class $y$ on one fibre, at a base point of a Weil family or
at a CM point of a Mumford family (hollow dots: the CM points, countable and
dense), is pushed into $W$ by the Gysin map and brought back to the invariant
summand by the base operator $\Lambda_{C}$. When $B(W)$ holds the result $u$
is an algebraic class on $W$, pinned down by the displayed pairing against the
invariant classes, and its restriction to a very general fibre is the
invariant part of $y$. The dashed arcs are isogenies between CM members along
Hecke orbits: they carry algebraic classes to algebraic classes but reach
only countably many points, and \Cref{lem:spreading}, in the box, needs
uncountably many. This is how \Cref{cor:weilequiv} and
\Cref{cor:mumfordequiv} turn each target into one case of \textup{(L)}.}
\label{fig:propagation}
\end{figure}

\begin{theorem}[The Mumford square at a CM point]\label{thm:mumfordcm}
Let $W\to C$ be a family as in \Cref{cor:mumfordlefschetz} and $c\in C$ a CM
point. Then every Hodge class on $X_{c}\times X_{c}$ is algebraic. In
particular the two exceptional classes of \Cref{prop:mumfordproduct} are
algebraic on $X_{c}\times X_{c}$.
\end{theorem}

\begin{proof}
The Hodge group $T$ of $X_{c}$ is a $\QQ$-torus contained in $G$. It lies in a
maximal $\QQ$-torus of $G=\Res_{F/\QQ}\SL_{1}(D)$, which is the restriction of
scalars of the norm one torus of a quadratic extension of $F$ inside $D$, of
rank three, one factor for each embedding of $F$. The cocharacter of the Hodge
structure lies in the factor of the embedding at which $D$ splits, and the
Galois group permutes the embeddings of $F$ transitively, so the smallest
$\QQ$-subtorus containing it has rank three: $T_{\CC}$ is a maximal torus of
$G_{\CC}\cong\SL_{2}^{3}$. Its eight weights on
$V=H^{1}(X_{c},\CC)=V_{1}\otimes V_{2}\otimes V_{3}$ are the vectors
$(\pm1,\pm1,\pm1)$, pairwise distinct. Hence
$\End^{0}(X_{c})\otimes\CC=\End_{T}(V)$ is the algebra of diagonal matrices
in a weight basis $(x_{w})$. For $n\ge1$, $\operatorname{Hom}(X_{c}^{n},X_{c})
\otimes\CC$ is the space of $n$-tuples $(\varphi_{1},\dots,\varphi_{n})$ of
such matrices, acting by $x_{w}\mapsto x_{w}\otimes u_{w}$ with
$u_{w}=(\varphi_{1,w},\dots,\varphi_{n,w})\in\CC^{n}$, so the vectors $u_{w}$
can be prescribed independently. The Hodge classes of $X_{c}^{n}$, tensored
with $\CC$, are spanned by the products of vectors $x_{w}\otimes u$,
$u\in\CC^{n}$, in $\bigwedge^{*}(V\otimes\CC^{n})$ in which the weights $w$,
counted with multiplicity, sum to zero.

We show that for every $n\ge1$ the ring $\Hdg^{*}(X_{c}^{n})$ is generated by
pull-backs of Hodge classes of $X_{c}$ along homomorphisms $X_{c}^{n}\to
X_{c}$. Put $w_{1}=(1,1,1)$, $w_{2}=(1,-1,-1)$, $w_{3}=(-1,1,-1)$ and
$w_{4}=(-1,-1,1)$; the weights are the $\pm w_{i}$,
$w_{1}+w_{2}+w_{3}+w_{4}=0$, and any three of the $w_{i}$ are linearly
independent, so the integer relations among them are the multiples of
$(1,1,1,1)$. Let $\alpha_{i}$ and $\beta_{i}$ be the multiplicities of $w_{i}$
and $-w_{i}$ in a nonempty multiset of weights with sum zero. Then
$\alpha_{i}-\beta_{i}$ is independent of $i$; so either some $\alpha_{i}$ and
$\beta_{i}$ are both positive, or all $\alpha_{i}$ are positive, or all
$\beta_{i}$ are. In each case the multiset contains one of the sets
$\{w_{i},-w_{i}\}$, $\{w_{1},w_{2},w_{3},w_{4}\}$,
$\{-w_{1},-w_{2},-w_{3},-w_{4}\}$, each of sum zero, and by induction every
product above is, up to sign, a product of monomials
$\bigwedge_{w\in P}(x_{w}\otimes u_{w})$ over such sets $P$, whose weights are
pairwise distinct. Choosing $f\in\Hom(X_{c}^{n},X_{c})\otimes\CC$ with
$f^{*}x_{w}=x_{w}\otimes u_{w}$ for $w\in P$ gives
$\bigwedge_{w\in P}(x_{w}\otimes u_{w})=f^{*}\bigl(\bigwedge_{w\in P}x_{w}\bigr)$,
with $\bigwedge_{w\in P}x_{w}$ a Hodge class of $X_{c}$ tensored with $\CC$.
The map $f\mapsto f^{*}(t)$ is polynomial and the rational homomorphisms are
Zariski dense in their complexification, so the same span is reached with $f$
defined over $\QQ$ and $t$ rational. The Hodge classes of the abelian fourfold
$X_{c}$ are algebraic because the Hodge conjecture holds for abelian fourfolds
\cite{Mar25a}, and pull-backs and products of algebraic classes are
algebraic. The case $n=2$ is the theorem.
\end{proof}

\begin{corollary}[The first classes beyond the Weil lines are one case of
\textup{(L)}]\label{cor:mumfordequiv}
Let $W\to C$ be as in \Cref{cor:mumfordlefschetz} and
$W^{(2)}=W\times_{C}W\to C$, an abelian scheme of relative dimension eight
with $W^{(2)}$ projective of dimension nine. The following are equivalent.
\begin{enumerate}[label=\textup{(\alph*)},leftmargin=2.6em,itemsep=2pt]
\item $B(W^{(2)})$ holds for the class $L_{W^{(2)}}$.
\item The two exceptional classes of $X_{t}\times X_{t}$ are algebraic for
  uncountably many $t\in C$.
\item The Hodge conjecture holds for $X_{t}\times X_{t}$ for every $t\in C$.
\item The Hodge conjecture holds for $W^{(2)}$; for this item $C$ is taken to
  be the compact Shimura curve of torsion-free level, or a finite \'etale
  cover of it.
\end{enumerate}
\end{corollary}

\begin{proof}
As in the proof of \Cref{cor:mumfordlefschetz}, the connected monodromy group
is $G$, so the monodromy invariants in $H^{*}(X_{t}\times X_{t},\QQ)$ are
contained in the $G$-invariants, which are the Hodge classes of
$X_{t}\times X_{t}$ for $t$ very general; conversely a $G$-invariant class is
flat, being fixed by the monodromy. The Hodge group of every fibre lies in
$G$, so the invariant classes are Hodge classes on every fibre. At a point $t$
that is not a CM point the Hodge group $H_{t}$ of $X_{t}$ is $G$. Indeed, if
$H_{t}$ were smaller, the Hodge classes of the powers of $X_{t}$ that are not
$G$-invariant would stay Hodge only on a discrete subset of the period
domain, the upper half plane; every point of the orbit of the Hodge structure
at $t$ under $H_{t}(\RR)^{+}$ carries them, so that orbit is a point, the
cocharacter of the Hodge structure is central in $H_{t}$, and $H_{t}$, the
smallest $\QQ$-group containing it, is a torus: $t$ would be a CM point. So
at such $t$ the Hodge classes of $X_{t}\times X_{t}$ are the invariant ones,
and $\End^{0}(X_{t})=\End_{G}(V)=\QQ$, because
$V\otimes\bar{\QQ}=V_{1}\otimes V_{2}\otimes V_{3}$ is an irreducible
representation of $G$; thus $X_{t}$ is as in \Cref{setup:mumford}.

(c) implies (b) trivially. (b) implies (c): the exceptional classes are
invariant, hence flat, so by \Cref{lem:spreading} a multiple of each is the
restriction of an algebraic class on $W^{(2)}$, and they are algebraic on
every fibre. At a point that is not a CM point, \Cref{thm:mumfordpowers} with
$n=2$ makes every Hodge class of $X_{t}\times X_{t}$ algebraic; at a CM point
this is \Cref{thm:mumfordcm}. (c) implies (a): by (c) every invariant class is
algebraic on every fibre, so by \Cref{lem:spreading} a multiple of it is the
restriction of an algebraic class on $W^{(2)}$, and
\Cref{thm:lefschetzfamily} gives $B(W^{(2)})$. (a) implies (c): at a CM point
$c_{0}$ every Hodge class of $X_{c_{0}}\times X_{c_{0}}$ is algebraic by
\Cref{thm:mumfordcm}, so by \Cref{thm:lefschetzpropagation}(iii) every
invariant class is the restriction of an algebraic class, and is algebraic on
every fibre. This is (c) at the points that are not CM points, where the
Hodge classes are the invariant ones; at a CM point (c) is
\Cref{thm:mumfordcm}.

For (d), split $H^{*}(W^{(2)},\QQ)$ into the summands $H^{p}(C,R^{q})$,
$p\in\{0,1,2\}$, as in the proof of \Cref{prop:mumfordfivefold}. The summands
with $p=1$ carry no Hodge class, by the argument given there, applied to the
fibre square. Those with $p=2$ are $f\cup H^{q}_{(q)}$, and
$f\cup u=i_{c*}(u|_{W^{(2)}_{c}})$ for any $c$; taking $c$ a CM point, where
$u|_{W^{(2)}_{c}}$ is a Hodge class and so algebraic by \Cref{thm:mumfordcm},
they are algebraic. The summands with $p=0$ are the invariants, all of Hodge
type. So (d) holds exactly when every invariant class is the restriction of an
algebraic class, which by \Cref{thm:lefschetzpropagation}(iii), applied at a
CM point, is (a).
\end{proof}

\begin{corollary}[One Weil family is one case of \textup{(L)}]
\label{cor:weilequiv}
Keep \Cref{setup:shimura}, let $s_{0}\in\cS$ be the image of a base point of
\Cref{thm:everyfamily}, and let $t_{1}\in\cS$ lie outside every proper
Zariski closed set $\Sigma_{\underline{P},\underline{q}}$ of
\Cref{thm:closedlocus}. Let $\kappa\colon C\to\cS$ be a morphism from a smooth
projective connected curve whose image contains $s_{0}$ and $t_{1}$, and
$W_{C}=\cA\times_{\cS}C$.
\begin{enumerate}[label=\textup{(\roman*)},leftmargin=2.6em,itemsep=2pt]
\item If $B(W_{C})$ holds for $L_{W_{C}}$, then $\Sigma=\cD_{n,n}$, and the
  Hodge conjecture holds in every degree for every member of the family with
  $\Hg(A)=\SU(V,H)$.
\item Conversely, if $\Sigma=\cD_{n,n}$ and the monodromy of $W_{C}$ is Zariski
  dense in $\SU(V,H)$, then $B(W_{C})$ holds for $L_{W_{C}}$.
\end{enumerate}
\end{corollary}

\begin{proof}
(i) The class $\kappa^{*}\omega$ is a global section of
$R^{2n}\pi_{*}\QQ$ on $C$, algebraic at a point over $s_{0}$ by
\Cref{thm:everyfamily}. By \Cref{thm:lefschetzpropagation}(ii) it is algebraic
at a point over $t_{1}$, so $t_{1}$ lies in some closed set
$\Sigma_{\underline{P},\underline{q}}$, which is therefore all of $\cS$.
\Cref{thm:closedlocus} gives $\Sigma=\cD_{n,n}$, and \Cref{thm:residual}(a)
the rest. (ii) With dense monodromy the invariants are the
$\SU(V,H)$-invariants, which are the Hodge classes of a member with
$\Hg(A)=\SU(V,H)$: the powers of the polarisation and the Weil classes, by
\Cref{thm:hdgcount}. All are algebraic on every fibre when
$\Sigma=\cD_{n,n}$, so they are restrictions of algebraic classes by
\Cref{lem:spreading}, and \Cref{thm:lefschetzfamily} applies.
\end{proof}

Curves $\kappa$ as in \Cref{cor:weilequiv} exist: the Baily--Borel
compactification of $\cS$ is projective \cite{BB66}, its boundary has
codimension $2n-1\ge3$, the rational boundary components of the domain of type
$\mathrm{I}_{n,n}$ being of type $\mathrm{I}_{n-r,n-r}$ with $r\ge1$, and a
general complete intersection curve through $s_{0}$ and $t_{1}$ avoids it.

\begin{theorem}[Two branches cannot be two objects]\label{thm:nosum}
Let $A$ be of $(K,-1,n)$-Weil type with $n\ge2$, $\omega\ne0$ a rational class
on the Weil line, and $E$ a perfect complex with $\ch(E)=N\omega$, $N\ne0$, and
$\sigma_{E}$ injective. Let
$E\cong\bigoplus_{i}F_{i}[k_{i}]$ be any decomposition. Then every summand
$F_{i}$ with $\ch(F_{i})\ne0$ has a Yoneda product
$\bigwedge^{2}\Ext^{1}(F_{i},F_{i})\to\Ext^{2}(F_{i},F_{i})$ that is not
injective. In particular $E$ is not a direct sum of shifts of line bundles,
of skyscraper sheaves of points, or of their images under autoequivalences of
$D^{b}(A)$, since the Ext algebra of each of these is the exterior algebra on
its $\Ext^{1}$.
\end{theorem}

\begin{proof}
By \Cref{thm:p2numerical}(ii) the kernel of $\operatorname{ev}_{E}$ in degree
two is $P\wedge Q$. The characteristic morphism is compatible with the
decomposition, so $\operatorname{ev}_{F_{i}}(p\wedge q)=
\operatorname{ev}_{F_{i}}(p)\operatorname{ev}_{F_{i}}(q)=0$ for all
$p\in P$ and $q\in Q$. Suppose the Yoneda product of $F=F_{i}$ is injective on
$\bigwedge^{2}\Ext^{1}$. Then the images $\bar P$ and $\bar Q$ of $P$ and $Q$
in $\Ext^{1}(F,F)$ satisfy $\bar p\wedge\bar q=0$ for all $\bar p,\bar q$, so
either $\bar P=0$, or $\bar Q=0$, or $\bar P=\bar Q$ is a line. In the first
case $\operatorname{ev}_{F}$ kills $P\wedge HH^{1}$, in the second
$Q\wedge HH^{1}$, in the third all of $HH^{2}=HH^{1}\wedge HH^{1}$; by the
compatibility $\sigma_{F}\circ\operatorname{ev}_{F}=(\,\cdot\,)\lrcorner
\ch(F)$ of \Cref{setup:hh} the same subspace annihilates $x=\ch(F)$.

These annihilators are computed in the free module of rank one that
$H^{*}(A,\CC)$ is over $HH^{*}(A)=\bigwedge^{*}(P\oplus Q)$, with generator the
holomorphic volume form $\Omega$: an element $y\cdot\Omega$ is killed by
$P\wedge HH^{1}$ exactly when $p\wedge y$ has top degree for every $p\in P$,
that is when $y\in\det P\wedge\bigwedge^{*}Q+\bigwedge^{\ge4n-1}(P\oplus Q)$.
Since $\det P\cdot\Omega=\pm\alpha_{+}$, the first summand gives the classes
$\bigwedge^{j}Q\cdot\alpha_{+}$, of type $(n-b,n+a)$ with $a$ wedges by
$V_{-}^{0,1}$ and $b$ contractions by $T_{+}$, and the second gives classes of
type $(1,2n)$, $(0,2n-1)$ and $(0,2n)$. The components of $x$ are of type
$(k,k)$, so $x\in\CC\alpha_{+}$. Symmetrically the second case gives
$x\in\CC\alpha_{-}$ and the third $x=0$. Finally $x$ is
rational, complex conjugation carries $V_{+}$ to $V_{-}$ and $H^{1,0}$ to
$H^{0,1}$, so $\bar\alpha_{+}\in\CC\alpha_{-}$, and a rational multiple of
$\alpha_{+}$ or of $\alpha_{-}$ is zero. So $\ch(F)=0$. As
$\sum_{i}(-1)^{k_{i}}\ch(F_{i})=N\omega\ne0$, some summand has nonzero
Chern character, and the first statement follows; the second is the special
case in which every summand has an exterior Ext algebra.
\end{proof}
